\documentclass[11.5pt,a4paper]{article}
\usepackage[utf8]{inputenc}
\usepackage[T1]{fontenc}
\usepackage[french]{babel}
\usepackage{geometry}
\geometry{top=2.5cm, bottom=2.5cm, left=2.5cm, right=2.5cm}

% --- Polices Modernes ---
\usepackage{helvet}
\renewcommand{\familydefault}{\sfdefault}

% --- Graphisme & Tableaux ---
\usepackage[dvipsnames]{xcolor}
\usepackage{tcolorbox}
\tcbuselibrary{skins, breakable}
\usepackage{array}
\usepackage{tabularx}
\usepackage{booktabs}

\definecolor{DeepMed}{HTML}{1B4F72}
\definecolor{Gold}{HTML}{B7950B}

% --- Styles de boîtes ---
\newtcolorbox{FocusScience}[1]{
  colback=blue!2,
  colframe=DeepMed,
  fonttitle=\bfseries,
  title=#1,
  boxrule=0.5pt,
  arc=1mm,
  breakable
}

\newtcolorbox{TableauSynthese}[1]{
  colback=white,
  colframe=gray!50,
  fonttitle=\bfseries\color{black},
  colbacktitle=gray!10,
  title=#1,
  arc=0mm,
  breakable
}

% --- Headers ---
\usepackage{fancyhdr}
\pagestyle{fancy}
\fancyhf{}
\lhead{\footnotesize Manifeste de l'Affirmation de Soi}
\rhead{\footnotesize Référence Scientifique \& Sociétale}
\rfoot{\thepage}

% --- Titre ---
\title{
    \vspace{-2cm}
    \Huge \textbf{L'Affirmation de Soi} \\
    \Large \textit{Neurosciences, Clinique et Perspectives Sociétales} \\
    \vspace{0.5cm}
    \large Document de Référence pour la Diffusion
}
\author{Antigravity Intelligence System -- Expertise Consolidée}
\date{\today}

\begin{document}

\maketitle

\begin{abstract}
Ce document synthétise soixante ans de recherche en psychologie comportementale et les dernières avancées en neuro-imagerie. Il explore le passage de l'inhibition (handicap du non-dit) à l'affirmation de soi, en détaillant les douleurs systémiques engendrées par la passivité et en proposant des clés de lecture neurologiques pour le futur de nos interactions sociales.
\end{abstract}

\tableofcontents
\newpage

\section{Fondements Scientifiques et Historiques}

\subsection{Évolution : De Wolpe aux Neurosciences (1960-2024)}
Le concept d'assertivité émerge en 1958 avec \textbf{Joseph Wolpe} comme une technique de désensibilisation face à l'anxiété. Depuis les années 1960, le mouvement humaniste l'a transformé en un levier d'épanouissement personnel. Aujourd'hui, grâce à l'IRM fonctionnelle (IRMf), nous ne parlons plus seulement de comportement, mais de structures cérébrales engagées.

\begin{FocusScience}{Focus Neurosciences : Ce que dit l'IRM}
L'imagerie par résonance magnétique a démontré que s'affirmer active le \textbf{Cortex Préfrontal Ventromédian (vmPFC)} et le \textbf{Striatum Ventral}. À l'inverse, l'inhibition (le fait d'oser ne pas dire) provoque une hyper-activité de l' \textbf{Amygdale} (peur) et une atrophie fonctionnelle de l' \textbf{Hippocampe} (mémoire), expliquant la confusion mentale lors des conflits.
\end{FocusScience}

\section{La Cartographie des Douleurs}

La non-affirmation de soi n'est pas une simple timidité ; c'est un état de stress chronique qui "use" l'organisme à travers tous les cycles de la vie.

\subsection{Tableau Synthétique : Situations et Douleurs Associées}

\begin{TableauSynthese}{Symptomatologie de l'Inhibition}
\begin{tabularx}{\textwidth}{|l|X|X|}
\hline
\textbf{Cycle de Vie} & \textbf{Situation de Blocage} & \textbf{Douleurs \& Symptômes} \\
\hline
\textbf{Enfance/École} & Harcèlement, peur du groupe. & Maux de ventre, phobie scolaire, repli. \\
\hline
\textbf{Adulte (Pro)} & Incapacité à poser des limites. & Burnout, épuisement, perte de sens. \\
\hline
\textbf{Adulte (Perso)} & Relation de couple asymétrique. & Ressentiment, frustration, perte d'identité. \\
\hline
\textbf{Séniors} & Perte d'autonomie non exprimée. & Dépression masquée, isolement, dénutrition. \\
\hline
\end{tabularx}
\end{TableauSynthese}

\section{Impact sur le Système Nerveux et Hormonal}

\subsection{La Cascade du Cortisol}
Le manque d'affirmation maintient le \textbf{Système Nerveux Sympathique} en état d'alerte. Cette activation prolongée entraîne une sécrétion excessive de cortisol. 
\begin{itemize}
    \item \textbf{Neurotoxicité} : L'excès de cortisol est toxique pour les neurones de l'hippocampe.
    \item \textbf{Manifestations Physiques} : Hypertension, troubles digestifs, inflammation chronique.
\end{itemize}

\section{Perspectives Cliniques et Psychiatriques}
Dans ses formes les plus extrêmes, la non-affirmation de soi s'apparente à des troubles diagnostiqués comme l' \textbf{Anxiété Sociale} ou le \textbf{Trouble de la Personnalité Évitante}. Le traitement ne peut se faire uniquement par la parole ; il nécessite une rééducation du "réflexe d'affirmation" pour recalibrer l'amygdale.

\section{Horizon Sociétal : Pourquoi est-ce le défi de demain ?}
Dans une société de plus en plus fragmentée et numérique, le pouvoir de dire "non" de manière constructive devient une \textit{soft skill} vitale. Sans affirmation de soi, l'individu se dissout dans les attendus sociaux, menant à une crise de santé mentale collective déjà visible via le burnout généralisé.

\section{Conclusion}
L'affirmation de soi est la brique de base du leadership authentique. Elle protège le cerveau du stress toxique et permet une interaction sociale durable. Ce document pose les fondations d'une éducation à l'assertivité nécessaire à tous les âges de la vie.

\end{document}
